\section{Holomorphic operator substitution}\label{sec:operator_substitution}

This section receives the operator-substitution result deliberately moved out of Paper IV. Fix an integer $N\geq1$ and the entire scalar function
\[
\Dfun_N(z)=\sum_{m=0}^{\infty}\frac{(-1)^m z^{Nm}}{(Nm)!}.
\]
More generally, the result applies to any specified holomorphic scalar function that is nonzero on the declared domain. On a simply connected zero-free domain $U$, choose $g=\log\Dfun_N$ once and define $\Sfun_N^{(r)}(z)=\exp(-r g(z))$ for complex $r$. The bounded Banach-space result below uses classical holomorphic functional calculus; it is separate from the finite canonical-jet calculus formalized in Part I.

\begin{standardfact}[Holomorphic functional calculus]\label{sf:holomorphic_calculus}
Let $A$ be a bounded operator on a complex Banach space.  If $f$ is holomorphic on a neighborhood of $\Spec(A)$, then
\[
f(A)
=
\frac1{2\pi i}\int_\Gamma f(z)(zI-A)^{-1}\,dz
\]
for an admissible positively oriented contour $\Gamma$ lying in the holomorphy domain and the resolvent set, with winding number one about the spectrum and zero outside the domain. The calculus preserves sums, products, and composition.  We use the standard form from \cite{DunfordSchwartz1958}; an accessible proof is given in \cite[Appendix A]{TaylorBanachAlgebras}.
\end{standardfact}

\begin{proposition}[Operator reciprocal family]\label{prop:operator_reciprocal}
Let $A$ be bounded and suppose $\Spec(A)$ lies in a simply connected region on which $\Dfun_N$ has no zeros.  Fix a holomorphic branch $g=\log\Dfun_N$ there.  Then
\[
\Sfun_N^{(r)}(A)
=
\exp\bigl(-r g(A)\bigr)
=
\bigl(\Dfun_N^{-r}\bigr)(A)
\]
is well defined and
\[
\Sfun_N^{(r)}(A)\Sfun_N^{(s)}(A)
=
\Sfun_N^{(r+s)}(A).
\]
If $(P_i)$ is a finite complete family of bounded idempotents, then
\[
P_i f(A)P_j
=
\frac1{2\pi i}
\int_\Gamma
f(z)P_i(zI-A)^{-1}P_j\,dz
\]
for every holomorphic $f$ on a neighborhood of the spectrum.
\end{proposition}

\begin{proof}
The spectral mapping theorem makes $\Dfun_N(A)$ invertible.  The product identity follows from multiplicativity of the functional calculus and the fixed logarithm.  Applying $P_i$ and $P_j$ to the Cauchy integral formula gives the packet formula.
\end{proof}

